\subsection{Landau--Lifshitz--Gilbert Equation}

\paragraph{Equation.}
The Gilbert form for the normalized magnetization is
\begin{equation}
    \frac{\dd\mvec}{\dd t}
    = -\gamma_0\,\mvec\times\Heff
      + \alpha\,\mvec\times\frac{\dd\mvec}{\dd t}.
    \label{eq:llg-gilbert}
\end{equation}
The equivalent explicit form is
\begin{equation}
    \frac{\dd\mvec}{\dd t}
    = -\frac{\gamma_0}{1+\alpha^2}
      \left[
        \mvec\times\Heff
        + \alpha\,\mvec\times(\mvec\times\Heff)
      \right].
    \label{eq:llg-explicit}
\end{equation}

\paragraph{Definitions.}
$\mvec$ is the normalized magnetization, $\gamma_0>0$ is the gyromagnetic factor,
$\alpha$ is the Gilbert damping parameter, and $\Heff$ is the effective field.

\paragraph{Assumptions.}
The magnetization magnitude is constant, so $\norm{\mvec}=1$.  Equation
\eqref{eq:llg-explicit} contains no additional spin torque or thermal field.

\paragraph{Units and sign convention.}
This form uses $\Heff$ in $\mathrm{A\,m^{-1}}$ and
$\gamma_0=\mu_0\abs{\gamma}$ in $\mathrm{m\,A^{-1}\,s^{-1}}$.  The electron
gyromagnetic sign is represented by the explicit minus sign.

\paragraph{Source.}
Add the paper or textbook used for the particular model here.

\paragraph{Implementation notes.}
Renormalizing $\mvec$ after a numerical step may hide integration error.  Prefer an
integration method designed to preserve the unit-vector constraint, and document any
renormalization explicitly.
